% Chapter 6 -- Rubric: PO11, PO9 (9.1--9.4), PO6 (6.1), PO7 (7.1), PO8 (8.1--8.4)
\chapter{Project Management, Teamwork, and Impact}
\label{ch:management}

\section{Project Management and Finance}

\subsection{Planning and Risk Management}
\label{sec:plan}

Five part-time members carrying a full course load (constraint C1) meant the
plan had to tolerate uneven availability rather than assume a steady rate of
work. It was therefore organised around \emph{work packages} closed by an
explicit verification gate rather than around calendar effort: a package was
counted as done not when its code existed but when it had been built,
integrated and verified against the running system. The version-control
history records exactly that discipline, in commits of the form ``built,
verified and integrated WP8''.

\begin{figure}[htbp]
  \centering
  \definecolor{gPrelim}{HTML}{6C7A99}
  \definecolor{gReq}{HTML}{4F8AA8}
  \definecolor{gBuild}{HTML}{2E8B7F}
  \definecolor{gRel}{HTML}{C08A33}
  \definecolor{gEval}{HTML}{8A5FA8}
  \resizebox{\textwidth}{!}{%
  \begin{ganttchart}[
      hgrid, vgrid, x unit=0.72cm, y unit chart=0.56cm, y unit title=0.56cm,
      time slot format=isodate-yearmonth, time slot unit=month,
      title label font=\scriptsize, bar label font=\scriptsize,
      group label font=\scriptsize\bfseries,
      bar height=0.6, bar top shift=0.2,
      bar/.style={draw=none, fill=gBuild},
      group/.style={draw=none, fill=black!55}
    ]{2025-09}{2026-09}
    \gantttitlecalendar{year, month=shortname} \\
    \ganttgroup{Preliminary study}{2025-09}{2025-12} \\
    \ganttbar[bar/.append style={fill=gPrelim}]{Problem selection with partner}{2025-09}{2025-10} \\
    \ganttbar[bar/.append style={fill=gPrelim}]{Background study of the field}{2025-10}{2025-12} \\
    \ganttbar[bar/.append style={fill=gReq}]{Requirements and feasibility}{2026-01}{2026-03} \\
    \ganttbar[bar/.append style={fill=gBuild}]{Prototype: planner + executor}{2026-03}{2026-04} \\
    \ganttbar[bar/.append style={fill=black!12}]{Examination period}{2026-05}{2026-05} \\
    \ganttbar[bar/.append style={fill=gBuild}]{Generalisation and architecture}{2026-06}{2026-07} \\
    \ganttbar[bar/.append style={fill=gBuild}]{Work packages WP0--WP9}{2026-07}{2026-08} \\
    \ganttbar[bar/.append style={fill=gRel}]{Reliability and fault attribution}{2026-08}{2026-09} \\
    \ganttbar[bar/.append style={fill=gRel}]{Web platform and second target}{2026-09}{2026-09} \\
    \ganttbar[bar/.append style={fill=gEval}]{Evaluation and report}{2026-08}{2026-09}
  \end{ganttchart}}
  \caption[Project schedule]{Project schedule, from problem selection through
  submission. Slate bars are the preliminary study that preceded any code; blue
  is requirements and feasibility; teal is construction; amber is the
  reliability work that had to precede measurement; violet is evaluation and
  writing. The pale bar in May is the examination period, which displaced
  milestone M2 in \Cref{tab:milestones}.}
  \label{fig:gantt}
\end{figure}

\paragraph{Phases.} The work fell into seven phases, each ending in a state the
next phase could build on. The first preceded the repository by six months and
produced no code, which is why it is easy to omit from a schedule and worth
recording here: the decisions it fixed constrained everything after it.

\begin{description}[leftmargin=2.2em,style=nextline]
  \item[P0. Preliminary study (September--December 2025).] Problem selection
    with the supervisor and the industry partner, and the background study of
    GUI test automation that established why the scripted and record-and-replay
    approaches in use do not survive interface change. This phase produced the
    project proposal and the reading that \Cref{ch:research} reports: it settled
    that the problem was worth attempting, that an agent-based approach was the
    live research direction, and that the measurement gap in the field
    (\Cref{subsec:measurement}) would shape what the project could claim. No
    code was written, and the architecture was not yet fixed.
  \item[P1. Requirements and feasibility (January--March 2026).] Analysis of the
    partner's technical requirements document ETA-TR-2026-001~\cite{etatr}, the
    literature review reported in \Cref{ch:research}, and the feasibility
    assessment in \Cref{sec:feasibility}. This phase fixed the four-process
    architecture and the decision to communicate only through a persistent
    graph.
  \item[P2. Prototype (March--April 2026).] The repository opened on 19 March 2026.
    A planner and an executor were integrated into a loop that ran end to end
    on a single hard-coded application, with the model provider treated as a
    swappable backend from the outset.
  \item[P3. Generalisation (June--July 2026).] Application-specific prompts,
    identifiers and ingestion paths were removed so that a new target arrives
    as configuration rather than as a code change. This is the phase that
    turned a demonstration into requirement DR-03, and it was the point at
    which the orchestration was rebuilt on a graph-execution framework.
  \item[P4. Work packages WP0--WP9 (July 2026).] The partner's eight requirement
    families and the team's derived requirements were decomposed into ten
    numbered packages, each built, integrated and verified in sequence, with
    the operator dashboard brought up alongside them.
  \item[P5. Reliability and fault attribution (August 2026).] Data-integrity
    defects, state-identity errors and loop instability were corrected;
    observability and a degradation register were added; the agent's own
    failures were made visible as a distinct class rather than being reported
    as application defects; and the executor and investigator roles were moved
    to a cheaper model tier once the quality cost of doing so had been
    measured.
  \item[P6. Second platform, evaluation and report (September 2026).] A browser
    executor brought the same planner and the same graph onto the web, a
    production commerce application and a live website were configured as
    targets, and the campaigns and writing reported in \Cref{ch:evaluation} and
    in this document were carried out.
\end{description}

\paragraph{Milestones, planned against actual.} \Cref{tab:milestones} records
the gates that closed each phase and where the outcome differed from the plan.

\begingroup
\small
\begin{xltabular}{\textwidth}{@{}>{\hsize=0.52\hsize}Y>{\hsize=0.26\hsize}Y>{\hsize=1.22\hsize}Y@{}}
  \caption{Milestones and their outcome against the plan}
  \label{tab:milestones}\\
  \toprule
  Milestone & Closed & Outcome against plan \\
  \midrule
  \endfirsthead
  \caption[]{Milestones and their outcome against the plan \emph{(continued)}}\\
  \toprule
  Milestone & Closed & Outcome against plan \\
  \midrule
  \endhead
  \midrule
  \multicolumn{3}{r@{}}{\footnotesize\itshape continued on the next page}\\
  \endfoot
  \bottomrule
  \endlastfoot

  M1 Loop runs end to end on one application & April 2026 & On plan. Planner and executor integrated on 7 April. \\
  \addlinespace
  M2 Application-agnostic pipeline & June 2026 & Late. Displaced by the examination period; recovered by compressing the phase into four weeks. \\
  \addlinespace
  M3 WP0--WP9 built and integrated & July 2026 & On plan, and the densest phase of the project. \\
  \addlinespace
  M4 Agent faults separated from application defects; cheaper model tier adopted & August 2026 & On plan. Both were prerequisites for any defensible measurement. \\
  \addlinespace
  M5 Second platform executing & September 2026 & Late relative to its value. Discussed below. \\
  \addlinespace
  M6 Evaluation campaigns and report & September 2026 & Compressed. The seeded-defect experiment identified in \Cref{ch:research} did not fit inside it and is carried into future work. \\
  \bottomrule
\end{xltabular}
\endgroup

\paragraph{Where the plan and the outcome differed.} Four deviations are worth
stating, because each of them changed what the project could claim.

\emph{May produced no work at all.} The repository records no commits between
28 April and 14 June. This was mid-session examinations, and it was foreseeable
rather than a surprise; the plan absorbed it by treating June as the
generalisation phase instead of spreading that work across May and June. The
cost was not lost capability but lost slack, and the slack would have been
useful in September.

\emph{The second platform arrived last.} Browser execution was built in the
final six weeks. Because the planner and the graph were already
application-agnostic by then, the browser executor could be added behind the
same interfaces in under two weeks, which is the strongest available evidence
that the generalisation phase paid for itself. The cost of the late arrival is
that the web target has carried far fewer campaign rounds than the Android
target, and the evaluation reports the two asymmetrically for that reason.

\emph{A finished component was deliberately left switched off.} A second,
tool-calling planner was built, and it is better on every property the team
cares about in principle: it grounds screen names against observed state,
rejects requirement identifiers that do not exist, and refuses duplicates
before committing them rather than after. It costs seven to eight model turns
against roughly four. It is \emph{not} the default, because it has not yet been
shown to produce better campaigns, only better-formed proposals. Adopting it on
the strength of its design would have been the easy decision; leaving it
switched off until a campaign settles the question is the correct one.

\emph{Evaluation was compressed.} Reliability work in August was necessary
before any measurement was meaningful, and it consumed time that had been
notionally reserved for campaigns. The consequence is visible in this report:
the evidence base is a set of campaigns rather than a standard benchmark suite,
and the open measurement identified in \Cref{subsec:measurement} remains open.

\paragraph{Risk register.} \Cref{tab:risks} lists the risks that were tracked,
how each one actually presented, and what was done about it. All but two
materialised, which is less a sign of good forecasting than of the fact that
most of them were characteristics of the problem rather than contingencies.

\begingroup
\small
\begin{xltabular}{\textwidth}{@{}>{\hsize=0.62\hsize}Y>{\hsize=0.78\hsize}Y>{\hsize=0.60\hsize}Y@{}}
  \caption{Risk register: exposure, mitigation and residual risk}
  \label{tab:risks}\\
  \toprule
  Risk & Mitigation applied & Residual \\
  \midrule
  \endfirsthead
  \caption[]{Risk register \emph{(continued)}}\\
  \toprule
  Risk & Mitigation applied & Residual \\
  \midrule
  \endhead
  \midrule
  \multicolumn{3}{r@{}}{\footnotesize\itshape continued on the next page}\\
  \endfoot
  \bottomrule
  \endlastfoot

  R1 The target exposes no stable identifiers and no text through the accessibility tree (C4). Materialised. &
    Screen identity resolved as a cascade of set-similarity measures with a perceptual image hash as the final tie-break, each threshold chosen analytically. &
    Real. Coverage on such applications remains below what a well-instrumented application yields, and \Cref{sec:equity} records this as an exclusion. \\
  \addlinespace
  R2 Model outputs are not reproducible across calls (C8). Materialised by definition. &
    Claims made over distributions of runs rather than single runs; the graph, not the model, holds state between rounds. &
    Intrinsic. It is a property of the method, not a defect to be removed. \\
  \addlinespace
  R3 A fixed prepaid inference balance is exhausted before the campaigns are run (C2). Did not materialise. &
    Cost per executed test made a design variable: the executor and investigator moved to a cheaper model tier after the quality cost was measured, and spend was reconciled against the provider's billing record. &
    Low. The published figure that DroidAgent could not repeat its own runs for cost reasons~\cite{droidagent} shows the risk was not hypothetical. \\
  \addlinespace
  R4 The agent damages a third party's live users or data (C5). Did not materialise. &
    Registration, one-time-password verification, payment and account deletion blocked at the configuration layer rather than discouraged (DR-07); disposable test accounts provisioned in advance. &
    Low, and deliberately bought at the price of coverage over exactly those flows. \\
  \addlinespace
  R5 No defect export or seeded-defect build from the partner (C3). Materialised. &
    Defect intelligence implemented and exercised against synthetic data; the limitation recorded rather than worked around. &
    High, and it is the binding constraint on the project's central claim. It is why precision and recall are not reported. \\
  \addlinespace
  R6 The agent's own failures are reported as application defects. Materialised, and was found in the record. &
    Three-way application/agent/environment attribution with exactly one definition in the source tree (NFR-D3); a degradation register that records every fallback taken (NFR-D2). &
    Moderate. Attribution is only as good as its classifier, and the classifier is itself a piece of the agent. \\
  \addlinespace
  R7 Concurrent runs corrupt the shared graph. Materialised. &
    A run lock admits one batch at a time, with a liveness check that detects a batch which never took a lock. &
    Low. \\
  \addlinespace
  R8 The loop stalls without terminating. Materialised repeatedly. &
    Step and round budgets, livelock detection, and the degradation register that made stalls visible instead of silent. &
    Moderate. Budgets bound the symptom; they do not remove every cause. \\
  \addlinespace
  R9 A human challenge, such as a CAPTCHA, blocks an unattended run. Materialised on the web target. &
    The run pauses, signals the operator audibly at intervals up to three minutes, and resumes on human input rather than failing the test or recording a false defect. &
    Accepted. It is a bounded loss of autonomy, and the alternative is a corrupted result. \\
  \addlinespace
  R10 A hosted model is deprecated or repriced. Did not materialise. &
    The provider is a configured backend, with three interchangeable options exercised during development. &
    Low in engineering terms; the commercial dependence itself remains (\Cref{sec:commercialisation}). \\
  \addlinespace
  R11 Uneven member availability under a full course load (C1). Materialised in May. &
    Work packages with verification gates, member branches merged through pull requests, and a written handoff document so that no package depended on one person being available. &
    Accepted. It cost the project its schedule slack. \\
\end{xltabular}
\endgroup

\subsection{Budgeting and Resource Identification}
\label{sec:budget}

The project had no procurement budget, no GPU allocation and no institutional
compute grant (constraint C2). That constraint determined the shape of the
entire budget, and it is worth being explicit about why. A system of this kind
needs a capable language model in the loop on every step; self-hosting one at
the sizes this system depends on requires accelerator hardware the team neither
had nor could obtain. Hosted inference through a commercial API aggregator was
therefore not a cost optimisation chosen over self-hosting; it was the only route available. The consequence is that the project's compute budget appears
entirely as metered API spend, with no hardware line at all, and that the
figure is small only because the team paid for tokens rather than for
accelerators.

Everything else the project needed (a machine to develop on, an environment to
execute against, an application with real behaviour and a specification to test
it against) was obtained without cash. \Cref{tab:budget} lists each
resource and how it was acquired. One convention governs the table and
determines its total: a resource the project did not pay for is recorded at
zero, not at an imputed market price. Assigning notional prices to
already-owned laptops and to the team's own labour would produce a larger and
more impressive figure, but not a truer one. Conversions use an approximate
rate of BDT 120 to the United States dollar.

\begingroup
\small
\begin{xltabular}{\textwidth}{@{}>{\hsize=0.58\hsize}Y>{\hsize=1.16\hsize}Y>{\hsize=0.26\hsize}Y@{}}
  \caption{Resources the project required, how each was obtained, and what it cost in cash. Inference is the only line that cost money, because it is the only input the team could not substitute with something already owned or freely available.}
  \label{tab:budget}\\
  \toprule
  Resource & How it was obtained & Cash cost (USD) \\
  \midrule
  \endfirsthead
  \caption[]{Resources, provenance and cash cost \emph{(continued)}}\\
  \toprule
  Resource & How it was obtained & Cash cost (USD) \\
  \midrule
  \endhead
  \midrule
  \multicolumn{3}{r@{}}{\footnotesize\itshape continued on the next page}\\
  \endfoot
  \bottomrule
  \endlastfoot

  Model inference, development & Metered credit on a commercial API aggregator, purchased out of pocket across the team's accounts over the project. This covers development, debugging, repeated specification re-ingestion and regression runs, and it is the dominant cash line: iteration on the pipeline costs far more inference than running the finished pipeline does. & $\approx$ 43 \\
  \addlinespace
  Model inference, evaluation campaigns & A separate prepaid project balance of USD 6.50 available at handover, of which USD 2.69 was consumed across 88 device executions. This is the only figure here read directly from a provider's billing record, and it is the basis of the per-test cost below. & 2.69 \\
  \addlinespace
  Android execution environment & Android emulators, running on each member's own laptop, carried the great majority of execution. This is what let five people develop and run tests in parallel without contending for a single handset, and it is also what made runs repeatable, since an emulator can be returned to a known state. One mid-range handset, already owned by a team member, covered the checks that needed real hardware, with the companion accessibility and input-method application installed (C6). & 0 \\
  \addlinespace
  Development hosts & Five personally owned laptops, each capable of running an emulator, the graph database and the four-process stack together. No server, and no accelerator of any kind. & 0 \\
  \addlinespace
  Application under test, and its specification & A production application with live users, and its requirements document, both supplied by the supervisor and the industry partner. & 0 \\
  \addlinespace
  Graph database, service and web frameworks, browser and device drivers, embedding and document-handling libraries & Open-source components, self-hosted, listed with their licences in \Cref{sec:ip}. No licence fee. & 0 \\
  \addlinespace
  Version control, continuous integration, observability, design export & Free tiers of hosted developer services. & 0 \\
  \addlinespace
  Engineering effort & Five part-time members over roughly six months. Not priced; see below. & \textit{n/a} \\
  \midrule
  \textbf{Total cash expenditure} & \textbf{Approximately BDT 5{,}500.} & $\approx$ \textbf{46} \\
  \bottomrule
\end{xltabular}
\endgroup

The development figure is reconstructed from the team's own account records
rather than consolidated from a single invoice, because members bought credit
separately; it is accurate to a few dollars and not to the cent. It also
excludes one member's spend, which was deliberately kept on smaller and cheaper
models and was not separately recorded, so the honest statement is that the
project's total metered inference cost was \textbf{under fifty United States
dollars}.

\paragraph{The input the table declines to price.} Engineering effort dominates
this project and is deliberately left uncosted. The team kept no timesheets, so
any hour count would be reconstructed after the fact; and any hourly rate
applied to it would be chosen rather than observed, with a factor of ten
between a student stipend and a loaded commercial salary. Multiplying one
invented number by another and printing the product as a total would make the
budget look complete at the cost of making it fictional. What can be said
without inventing anything is the shape: the project's cost is almost entirely
labour that the course absorbed, and its cash cost is under fifty dollars.
\Cref{sec:economics} takes up what the same system would cost an organisation
that had to pay for both.

\paragraph{What the metered part actually costs.} Because inference is the only
input that scales with use, it is worth separating into its two very different
components, both measured rather than derived from a price card. Executing one
test costs about USD 0.008, roughly one taka: built from 10.1 device steps per
test across 232 stored observations, an accessibility tree averaging 4,686
bytes per step, and the planner's 384{,}000 logged tokens over 108 calls. Within
that figure the device agent accounts for about 92\% and the planner about 8\%,
which is why reducing median steps per test was a real saving and why
optimising the planner's prompt would not have been. A twenty-five-test
campaign is therefore about USD 0.20, and a five-hundred-test nightly suite
about USD 4.

Ingesting a specification behaves in the opposite way. It costs about USD 1.80,
the price of roughly two hundred test executions, because extraction
runs on the most expensive model tier, splits the document into sections and
issues one uncapped call per section, three times over for self-consistency. It
is paid once per document rather than once per run, so it is a fixed cost in
the terms of \Cref{sec:economics} and not an operating one.

Those two figures explain the gap between the two inference lines in
\Cref{tab:budget}. Running the finished system is cheap; building it was not,
because development re-ran the expensive half repeatedly: every change to the ingestion pipeline meant re-reading the specification from scratch, and the
executor and investigator ran on a more expensive tier until the quality cost
of moving them down had been measured. The metering controls that would close that gap (an output ceiling on extraction, a lower sample count once it has
stabilised, and a cache that refuses to re-read an unchanged document) are identified in \Cref{sec:conclusion} as outstanding work. They matter more for a
team iterating on the pipeline than for an organisation running it.

\subsection{Projected Cost of Self-Hosting}
\label{sec:selfhost}

Nothing in this subsection was spent. It is a projection, kept separate from
\Cref{tab:budget} for that reason, and it exists because the budget's headline
figure is only honest if the alternative it avoided is priced. The project
rented inference because it had no accelerator (C2); a reader entitled to ask
what the system would cost without that dependency deserves an answer with its
arithmetic shown rather than an assurance.

\paragraph{What would actually have to be served.} Not every model in
\Cref{tab:models} matters here. The device agent dominates: it is a vision
model invoked on every step, it accounts for about 92\% of the cost of an
executed test, and its published checkpoint is a 32-billion-parameter one.
Sizing follows from that count. Weights at 16-bit precision need roughly
64\,GB; at 8-bit roughly 32\,GB; at 4-bit roughly 16--18\,GB, before the
key-value cache, the vision encoder's activations and one
$1080\times1920$ screenshot per step. The practical minimum is therefore a
single 48\,GB accelerator running an 8-bit quantisation, or two 24\,GB consumer
cards running a 4-bit one. The planner and the investigator are smaller and
would fit alongside; specification extraction is the one workload that would
still argue for a hosted call, since it runs on a far larger model but only
once per document.

\paragraph{The arithmetic.} \Cref{tab:selfhost} sets the two routes against the
measured API cost of USD 0.008 per test. Prices are indicative at the time of
writing and are the input a reader should replace with their own; the structure
of the comparison is the part that does not change.

\begingroup
\small
\begin{xltabular}{\textwidth}{@{}>{\hsize=0.70\hsize}Y>{\hsize=1.34\hsize}Y>{\hsize=0.96\hsize}Y@{}}
  \caption{Projected cost of serving the device agent on owned or rented hardware, against the measured API cost. Utilisation, not price, is what decides the comparison.}
  \label{tab:selfhost}\\
  \toprule
  Route & Basis & Cost per test \\
  \midrule
  \endfirsthead
  \caption[]{Projected cost of self-hosting \emph{(continued)}}\\
  \toprule
  Route & Basis & Cost per test \\
  \midrule
  \endhead
  \midrule
  \multicolumn{3}{r@{}}{\footnotesize\itshape continued on the next page}\\
  \endfoot
  \bottomrule
  \endlastfoot

  Hosted API (what the project did) & Measured against the provider's billing record, at 10.1 device steps per test. & USD 0.008 \\
  \addlinespace
  Rented accelerator & One 48\,GB cloud instance at roughly USD 1.25 per hour, run continuously for a year: about USD 11{,}000. At one test every two minutes of occupied card, about 260{,}000 tests. & $\approx$ USD 0.042 \\
  \addlinespace
  Purchased accelerator & A 48\,GB workstation card at USD 5{,}000--8{,}000 in a host costing USD 8{,}000--12{,}000 all in, amortised over three years and ignoring power, cooling and administration. & $\approx$ USD 0.013 \\
  \bottomrule
\end{xltabular}
\endgroup

\paragraph{Why renting a card costs more than renting tokens.} The result is
counter-intuitive only until the duty cycle is examined. A test occupies its
device for up to the 900-second wall-clock bound configured for these
campaigns, but consumes accelerator time only in the brief interval between
observing a screen and deciding the next action, on the order of a hundred
seconds of inference across 10.1 steps. The rest is the emulator settling,
animations completing and the application responding. A dedicated card
therefore idles for the great majority of a run, and the tenant pays for that
idling; a hosted provider does not charge for it, because it fills those gaps
with other tenants' work. Per-token billing is, in this specific sense, a
better match to the workload than per-hour billing.

That also identifies the condition under which the conclusion reverses. The
card becomes economical exactly when it can be kept busy: many emulators driven in parallel against one served model, which is the natural shape of a
nightly regression fleet and not the shape of this project. A purchased card at
high utilisation approaches the API price and would undercut it at scale; the
same card at this project's utilisation would not.

\paragraph{The argument that is not about money.} Cost is the weaker half of
the case for self-hosting. The stronger half is that every screenshot of the
application under test is currently transmitted to a third party, and the
application under test is a production system with live users. The
configuration boundary of \Cref{sec:ethics} keeps the agent away from
destructive actions, but it does not stop interface contents leaving the
machine. An organisation with a data-residency obligation, or one testing an
unreleased product, would have to self-host regardless of what
\Cref{tab:selfhost} says, and because the models in \Cref{tab:models} are
open-weight checkpoints rather than proprietary endpoints, that path remains
open to them. Retaining that option was the reason the model provider was kept
behind a swappable interface from the first prototype.

\subsection{Economic Analysis and Financial Prospecting}
\label{sec:economics}

\paragraph{What it cost to build.} Under USD 50 in cash
(\Cref{tab:budget}), essentially all of it metered model inference, plus a
session of part-time work by five students. The cash figure is the one worth
quoting and it is the strongest fact in this section: an autonomous exploratory
testing agent was built, debugged and evaluated for under fifty dollars, by a
team with no GPU allocation at all. The expensive input to this class of system
is inference, and renting inference by the token has become cheap enough to be
a rounding error rather than a barrier to entry, which is precisely what makes the approach reachable for a group that could not have afforded the
hardware to self-host. The labour is deliberately left unpriced in
\Cref{sec:budget}, and the reason carries over here. A commercial team would
not reproduce this cost structure at all: five part-time students over a
session is not two full-time engineers over three months, and a firm paying
market salaries with overhead should expect a first production-grade version of
this system to cost an order of magnitude more than anything this project can
evidence. What transfers to such a firm is not the total. It is the unit
economics below.

\paragraph{What it costs to run.} Operating cost has three components. The
metered one is model inference at an order of one United States cent per
executed test case, derived in \Cref{sec:budget}, reported in full in
\Cref{ch:evaluation} and
set against published figures of USD 13--22 per application for a two-hour
budget~\cite{droidagent} and USD 10.17--18.76 per thousand
queries~\cite{autodroid} in \Cref{sec:feasibility}. The second is
infrastructure: a single host running four processes, a graph database and
either a device or a browser, which is a developer workstation rather than a
cluster. The third is the one that actually dominates and is the least
discussed in this literature: \textbf{a human being must triage every candidate
defect the agent reports.} The system is scoped so that a reported defect is a
candidate requiring confirmation, and that confirmation is paid for in
engineer-minutes at full rate.

\paragraph{Where the value is.} The value displaced is manual exploratory
testing and the maintenance of scripted suites, for the three user groups
identified in \Cref{sec:market}. One published deployment gives an order of
magnitude for the saving: AUITestAgent's authors report 1.5 person-days saved
per regression round on a \emph{single} business line at
Meituan~\cite{auitestagent}. Valued at a junior engineer's cost in this market,
1.5 person-days is a few tens of dollars; valued at a mid-level engineer's
fully loaded cost in a market like Singapore or Seoul it is several hundred.
On a fortnightly release cadence, a single business line therefore sits
somewhere between USD 900 and several thousand dollars of displaced effort per
year. The inference required to displace it, at one cent per test, is tens of
dollars annually even for a suite of several hundred tests per round. The ratio is
favourable by two to three orders of magnitude, and it would remain favourable
under fairly violent revisions of the assumptions.

\paragraph{The number that decides it, and that this project cannot supply.}
The preceding paragraph is an argument about cost, and cost is the easy half.
The value of the output depends on \emph{precision}: every false defect report
consumes an engineer's investigation time, so a tool reporting candidates at
low precision transfers cost to developers rather than removing it, and does so
invisibly. This project cannot report precision or recall, because the partner
supplied neither a defect export nor a seeded-defect build (constraint C3), and
because, as \Cref{subsec:measurement} establishes, none of the seven systems
reviewed reports them either. The correct reading is therefore: the cost side
of the business case is established and unusually strong, and the value side is
bounded by a measurement the whole field is missing. A prospective adopter
should treat the seeded-defect experiment named in \Cref{sec:conclusion} as the
condition on any return-on-investment claim, this one included.

\paragraph{Revenue models, if it were taken further.} Three are plausible, in
descending order of fit. As an \emph{internal tool} at an organisation that
already runs a QA function, the return is displaced manual effort and no
revenue model is needed; this is the partner's own framing and the most likely
outcome. As a \emph{metered service}, billing per executed test or per campaign
aligns price with the cost driver and suits the smaller firm in
\Cref{sec:market} that cannot staff a QA team at all, for whom the alternative
is not a cheaper tool but no systematic testing. As \emph{open core}, the agent
and graph are released and the revenue comes from hosting, target onboarding
and support, which is the model best matched to a system whose hard part is
configuring a new application rather than running it. Per-seat licensing fits
worst, because the system's cost scales with tests executed and not with the
number of people watching.

\subsection{Sustainability in Business and Commercialisation}
\label{sec:commercialisation}

\paragraph{The path that exists.} This system was not built speculatively. It
answers a technical requirements document written by Samsung Research and
Development Institute Bangladesh, which scoped an eight-family enhancement
programme at 17--23 weeks of effort with a phased dependency
ordering~\cite{etatr}. All eight families were delivered
(\Cref{tab:partner-req}). The natural commercialisation path is therefore
handover rather than launch: the partner already has the QA function, the
release cadence and the applications that make the system worth operating, and
it specified the capability itself. What this project owes that path is a
system that can be picked up, which is why the operator surface is a dashboard
and a configured JSON target profile rather than a set of scripts (DR-09), and
why a campaign exports as a PDF a non-participant can read (DR-08).

\paragraph{What stands between the current system and that handover.} Three
things, stated without softening. The deployment model is the first and
largest: the system is scoped to run as a local, developer-operated stack on a
single trusted host (constraint C7), and moving it into shared infrastructure
requires authentication between the four processes, secret management that does
not rely on a local environment file, multi-tenant separation of project-scoped
graphs, and an integration into continuous integration that this project did
not attempt. The second is the measurement gap of \Cref{sec:economics}: an
organisation should want a precision figure before it routes agent output into
a developer's queue. The third is device supply, which is mundane and real:
unattended Android execution needs a device or emulator with the companion
application installed, and an organisation testing across many device profiles
needs a device farm, which changes the operating cost from the figure in
\Cref{tab:budget} to something considerably larger.

\paragraph{What makes it viable to maintain.} The maintenance surface was kept
deliberately small. A new application is onboarded as a validated configuration
profile rather than a code change, so the recurring work of adoption falls on
an operator rather than on a developer. The model provider is a configured
backend with interchangeable options, so a deprecation or a repricing is a
configuration change and not a migration, which matters because the single
largest external dependency of this system is a commercial inference API whose
terms the project does not control. The knowledge graph is the only durable
state, and it has a backup and restore path. The component most likely to need
continuing work is the one that is least general: the element-location and
screen-identity strategy, which had to be extended for an application exposing
no stable identifiers (constraint C4) and will need extending again for the
next toolkit that exposes something different.

\paragraph{Who would operate it.} Within the partner, a QA engineer rather than
the authors: the system is operable by one person from the dashboard, which was
a requirement and not an accident. Outside it, the realistic route is the
open-core one in \Cref{sec:economics}, with the caveat recorded in
\Cref{sec:ip} that the graph database is distributed under a copyleft licence,
which shapes how the system may be redistributed even though it does not
constrain its use here. Beyond the course, the team's own position is stated in
\Cref{sec:conclusion}: the most valuable next contribution is not a feature but
the seeded-defect experiment, and that is work an academic group is better
placed to do than a product team.

\section{Individual and Teamwork}
\label{sec:team}

\subsection{Participation and Contribution}
\label{sec:participation}

\Cref{tab:contrib} states each member's role and principal contributions. The
allocation was by subsystem rather than by task, so that one person carried a
component from design through to its failures, which is the arrangement that
makes the accountability statement in \Cref{sec:accountability} meaningful.

One caution about how this table should be read. It is reconstructed from the
version-control record, and commit counts are a poor proxy for contribution:
they reward frequent small commits, they are silent about design work,
literature review, review of other people's code and debugging that ends in no
change, and they understate work done in pairs. The table therefore names
\emph{what each member owned and delivered} and does not rank the members.

\begingroup
\small
\begin{xltabular}{\textwidth}{@{}>{\hsize=0.44\hsize}Y>{\hsize=0.32\hsize}Y>{\hsize=1.52\hsize}Y>{\hsize=0.22\hsize}Y@{}}
  \caption{Member roles and principal contributions, in the order the team lists itself. The final column gives the month of that member's first commit in the version-control record.}
  \label{tab:contrib}\\
  \toprule
  Member & Role & Principal contributions & Active from \\
  \midrule
  \endfirsthead
  \caption[]{Member roles and principal contributions \emph{(continued)}}\\
  \toprule
  Member & Role & Principal contributions & Active from \\
  \midrule
  \endhead
  \midrule
  \multicolumn{4}{r@{}}{\footnotesize\itshape continued on the next page}\\
  \endfoot
  \bottomrule
  \endlastfoot

  Mohammad Raihan Rashid (2105046) &
    Project lead; architecture, ingestion, dashboard &
    Opened the repository on 19 March 2026 and fixed the four-process architecture and the decision that the agents communicate only through a persistent graph. Owns the ingestion pipeline: removed the hard-coded document path and made ingestion modular, which is what lets a new application arrive as configuration rather than as a code change (DR-03). Built the operator dashboard and carried it through its successive revisions. Substantial planner work from the first prototype onward, including the early model-backend experiments that made the provider swappable. Centralised configuration and fixed the state-identity and observation defects underlying screen recognition. Made the agent's own failures visible as a distinct class (NFR-D3), which is the precondition for any defensible measurement. Wrote the work plans: decomposed the partner's requirement families into the numbered work packages, defined the verification gate for each, and built WP0--WP3 before handing the sequence on. Verified completed packages against the running system. Reviewed and merged the team's pull requests throughout, and drafted the greater part of this report. &
    Mar 2026 \\
  \addlinespace
  Mehemud Azad (2105014) &
    Planner and retrieval &
    Integrated the executor with the planner to close the first end-to-end loop. Built the goal-based planner and the trajectory evaluator, and contributed the largest share of commits to the planner package. Owns the retrieval service and the embedding backend, and rebuilt the orchestration on a graph-execution framework. Added the observability layer, carried out the migration to a cheaper model tier, and added graph backup and restore. Contributed the largest share of the project's internal documentation. &
    Apr 2026 \\
  \addlinespace
  Jonayed Mohiuddin Shanto (2105060) &
    Execution and integration &
    Generalised the device gateway by removing application-specific prompts and identifiers, which opened the path to DR-03. Took the work-package sequence from WP1 onward: verified WP1--WP3 and filled their gaps, implemented WP4, and built, verified and integrated WP5--WP9, the largest single body of delivery in the record. Wrote the browser executor from scratch, bringing the system onto its second platform, and authored the target profiles for both platforms. Fixed the loop-stability, concurrency and run-lock defects, and added the human-in-the-loop CAPTCHA pause. Reclassified a dead-browser condition that had been recorded as an agent navigation fault, a correction to the attribution taxonomy itself. Wrote \Cref{subsec:formative}, which traces the planner, player and investigator decomposition back to the multi-agent structure it was adopted from, and states what the project took from that work and what it deliberately did not. &
    Apr 2026 \\
  \addlinespace
  Khalid Hasan Tuhin (2105002) &
    Operator surface and web target &
    Built the target-configuration interface and the knowledge-upload path in the dashboard, which is what makes DR-09 an operator action rather than a code change. Contributed to the browser executor: the wandering guard, risk-priority wiring, agent-difficulty steering and semantic retrieval of execution notes, followed by a reliability pass with a regression suite. Configured the production commerce target and the live web target. Fixed the executor's crash on provider rate-limiting and a precondition misclassification. &
    Jun 2026 \\
  \addlinespace
  Niloy Kumar Mondal (2105044) &
    Run reporting, literature grounding, internal review &
    Built the run-report agent that turns a campaign into a document a stakeholder can act on: it gathers execution logs, findings and coverage from the graph, filters outcomes against the fault-attribution taxonomy so that only application-caused failures are reported as defects, and renders a PDF with a generated narrative (DR-08). Surveyed fifteen recent related works, the material behind \Cref{subsec:additional}. Carried out the structured internal code review against an explicit severity rubric referred to in \Cref{sec:accountability}, and ran testing against the web target. &
    Sep 2026 \\
\end{xltabular}
\endgroup

\paragraph{Participation in decisions.} Design decisions were taken in
discussion and recorded where they had consequences: the choice to communicate
only through the graph, the decision to fix a separate fault-attribution
taxonomy rather than let agent failures be reported as defects, the migration
to a cheaper model tier, and the decision not to promote the tool-calling
planner to default. Each of these was argued by more than one member, and the
last two were settled against the position initially preferred by the person
who had done the work, which is the more informative fact about how the team
operated.

\guidance{TODO (team): each member should read their own row and correct it
before submission, and add any contribution the version-control record does not show, in particular design discussion, review of others' code, and debugging
that ended in no commit. Schedule slips are recorded against milestones in
\Cref{tab:milestones} rather than per member; if any individual responsibility
was not met on time, add it there with the reason.}

\subsection{Collaboration and Conflict Resolution}
\label{sec:collaboration}

The team worked on member branches merged into the mainline through pull
requests reviewed by the repository owner. This was not ceremony: with five
people editing a planner, a retrieval service and a shared graph schema, the
alternative to review was schema drift, and schema drift is precisely the class
of defect that appeared when it was not followed and had to be repaired in
August. Written artefacts carried what conversation could not: a workflow
document, a planner-prompt anatomy, a reliability note and a handoff roadmap,
so that a member returning after an absence could resume without needing
someone else to be available.

Disagreements were technical, and the useful thing to record is the mechanism
that settled them, because the mechanism was consistent.

\emph{Disagreement settled by measurement, not by authorship.} The tool-calling
planner is better designed than the default one on every property the team
values, and its author argued for adopting it. The team did not, because being
better formed is not the same as producing better campaigns, and no campaign
had yet been run to decide it. Both planners remain in the source tree behind a
configuration switch. Leaving a colleague's finished work switched off is the
harder outcome socially and the correct one methodologically, and it was
reached by agreeing in advance what evidence would decide the question.

\emph{Disagreement settled by looking at the record.} Four executions had been
recorded as navigation failures, which is an agent fault and implies the
team's own navigation logic was broken. Inspection showed the browser had
simply ceased to exist, and the condition was reclassified as an environment
fault that ends the batch immediately. The general lesson the team took from it
is the one that governs the whole project: when a result is ambiguous between
``our system is wrong'' and ``the environment was wrong'', the disagreement is
not resolved by argument but by making the system record which one it was.

\emph{Disagreement anticipated by process.} The structured internal code review
against an explicit severity rubric was adopted so that criticism of code
attached to a written standard rather than to a person, and so that the
resulting work could be scheduled rather than negotiated.

\subsection{Leadership and Direction}
\label{sec:leadership}

Direction was distributed by subsystem rather than concentrated in one member.
Each of the five owned a component and was the person answerable for it, and
the coordinating decisions, which sequence of work packages to build, when a
package was allowed to count as done, and what the team would stop doing, were
taken jointly. The repository owner held the merge decision, which is the one
piece of authority that cannot be shared without the shared graph schema
decaying.

Three devices supplied direction in practice. The first is the work-package
decomposition, which converted an open-ended research brief and the partner's
phased dependency ordering~\cite{etatr} into ten numbered units with a
verification gate each, so that progress was legible to everyone rather than
resident in the person doing the work. The second is a written handoff
document, produced in September, that states where the project stands, what the
agent has actually found, and what to do next \emph{in priority order},
written on the assumption that the reader is new, which is the form that
survives a member becoming unavailable. The third is the supervisor's steering,
which supplied a production application with live users as a second target; the
effect of accepting it was to force the generality requirement and the safety
boundary from aspirations into enforced constraints, and it was the single
most direction-changing input the project received from outside itself.

Setbacks were handled by instrumenting rather than by exhorting. When the loop
stalled intermittently, the response was livelock detection and a degradation
register that made stalls visible. When data-integrity defects were found, the
response was a regression suite that now runs six modules and just over a
hundred checks. When the agent's failures were found to be masquerading as
application defects, the response was to give fault attribution exactly one
definition in the source tree. The pattern is deliberate: a team of five
part-time members cannot rely on everyone remembering a caveat, so a caveat that
matters was made into a check that fails.

Members' views were listened to in the specific sense that the two most
consequential decisions of the late project, the model-tier migration and the
refusal to promote the tool planner, both went against the initial preference of
the member who had done the work in question, and both were accepted by that
member because the criterion had been agreed beforehand. Motivation over a long
session came less from encouragement than from the campaigns: an agent that
proposes its own test, drives an application nobody wrote a script for, and
returns a finding a human had not anticipated is a persuasive thing to watch,
and running one was the reliable way to restore momentum after a slow month.

\subsection{Multidisciplinary Engagement}
\label{sec:multidisciplinary}

The project is a computer-science project, but four of its binding decisions
were settled outside computer science, and the team had to engage with each
field on its own terms rather than through a software proxy.

\paragraph{Industry requirements engineering.} The starting point was a
partner's technical requirements document, not a research question the team
chose~\cite{etatr}. Working from it meant accepting an external party's
decomposition, effort estimate and dependency ordering, reconciling eight
requirement families against what was technically reachable in a session, and
reporting delivery in the partner's own identifiers so that the result could be
audited against the commission. The traceability matrix in
\Cref{tab:partner-req} is the artefact of that engagement, and the ten
requirements the team derived beyond it (\Cref{tab:derived-req}) are what
contact with a real application added to a specification written in advance.

The engagement was iterative rather than a single handover. An early iteration
of the system was submitted to the partner for review, and the criticism that
came back was not about missing features but about generality: the ingestion
path was fixed, and application-specific values were hard-coded through the
pipeline, so what had been demonstrated on one application would not move to
another without being rewritten. The team accepted that assessment and rebuilt
the affected parts in June, removing the hard-coded paths and identifiers and
making ingestion modular. Requirement DR-03 and phase P3 of \Cref{fig:gantt}
are both consequences of that review. Taking an adverse assessment from an
industrial reviewer and reworking already-working code in response is a
different discipline from writing the code in the first place, and it is the
part of the partnership that changed the architecture most.

\paragraph{Applied domains with live users.} The principal evaluation target
is Shobar Khamar, a production agricultural-commerce application with real farmers and sellers
on it, and a payment and one-time-password flow. Testing it required
understanding its domain workflow well enough to judge what correct behaviour
looked like, and required a judgement that is not a software judgement at all:
which operations may never be exercised against other people's livelihoods.
That analysis produced constraint C5 and requirement DR-07, and it is discussed
as an ethical obligation rather than a technical one in \Cref{sec:codes}.

Commerce was not the only domain the system was pointed at. Once the browser
executor existed, the same planner was run against a range of other web
applications, among them DataGhurhi, a departmental form and survey platform
whose submission workflows share almost no interaction vocabulary with a
marketplace. The breadth mattered for a specific reason: correct behaviour in a
survey builder and correct behaviour in a checkout flow have nothing in common,
so each target had to be understood in its own terms before the agent's output
could be judged at all. A system that cannot be told what the application is
for will report the difference between domains as though it were a defect.

\paragraph{Standards, law and regulation.} \Cref{sec:standards} works through
software-testing and quality standards, requirements-engineering practice and
the OWASP guidance for language-model applications, and applies them as
compliance obligations with stated controls: application content treated as
untrusted data in every prompt, credentials kept out of the planner path
entirely, and destructive actions blocked at the configuration layer. Reading
an instrument to determine what it requires of you is a different activity from
engineering, and the project's safety position rests on it.

\paragraph{Human factors and the economics of the output.} The design of what
the system reports was settled by evidence about people rather than by
preference: the developer study cited in \Cref{sec:market} finds that the
preferred form of an automatically generated test is natural language carrying
the expected output, the steps and the feature under test, and the system's
output takes that form for that reason. The market, feasibility and economic
analysis in \Cref{sec:market}, \Cref{sec:feasibility} and \Cref{sec:economics}
is business analysis, and it is what establishes that the interesting question
about this system is not whether it can be built but whether its reports are
precise enough to be worth a developer's time.

\section{Impact on Society}
\label{sec:society}

This system does not reach the public directly. It tests other people's
software, so its effect on society travels through two channels: the quality of
the applications that reach users, and the work of the people who currently
test them. Both are considered below, followed by the negative impacts and what
the design does about them.

\paragraph{Software quality, and why it is a safety question here.} A defect in
the graphical interface of a consumer application is usually an annoyance. In
the class of application this project was evaluated against it is not. The
target is an agricultural commerce platform used by smallholder farmers and
buyers in Bangladesh, where a wrong price, an order routed to the wrong seller,
or a total that does not match its line items is a direct financial loss to a
party with very little margin to absorb it. Those are exactly the defects that
a crash detector cannot see, for the reason set out in
\Cref{subsec:measurement}, and exactly the class that a specification-grounded
oracle is built to catch. \emph{Who is affected:} the end users of any
application tested with this system. \emph{How the design responds:} the oracle
is derived from the requirements rather than from the absence of a crash, and
every generated test cites the requirement it covers, so a defect report can be
traced to the behaviour that was promised.

\paragraph{Access to systematic testing.} Cost is not a neutral engineering
detail in this context. A firm that cannot fund a dedicated quality-assurance
team does not substitute a cheaper tool; it ships without systematic testing at
all. Measured cost per executed test is therefore a societal variable as well
as a budgetary one, and it is the reason NFR-D1 exists.
\emph{Who is affected:} small development firms, disproportionately in
economies like Bangladesh's, and the users of the software they ship.
\emph{How the design responds:} unit cost is measured rather than assumed
(DR-10), the expensive roles were moved to cheaper model tiers once the quality
cost of doing so had been measured, and a new target is configured as a
validated JSON profile through the dashboard rather than as a code change
(DR-09), so adoption does not require an engineer who can modify the system.

\paragraph{Accessibility, as a by-product.} The Android executor drives
applications through the platform accessibility tree and the web executor
through the accessibility-annotated DOM. The consequence is worth stating
plainly: a control this agent cannot locate is, in general, a control that a
screen-reader user cannot locate either. The agent's element-location failures
are therefore an accessibility signal, and its failure log doubles as a rough
accessibility audit of the application under test.
\emph{Who is affected:} users of assistive technology.
\emph{How the design responds:} controls are addressed by accessible role, name and content description rather than by pixel coordinates (\Cref{sec:standards}), and each agent-attributed failure records the specific control that could not be reached. The agent does also receive a screenshot, which sharpens the signal rather than weakening it: a control the agent can see in the image but cannot name from the accessibility tree is exactly a control that is visible to a sighted user and invisible to a screen reader, and that divergence is the accessibility gap itself. This is a genuine benefit but an incidental
one; it should not be presented as a designed accessibility audit, because the
agent does not check contrast, focus order, or any other WCAG criterion.

\paragraph{Language, and a failure that was silent.} The project's text
normalisation originally filtered input to the Latin range, which strips every
non-Latin character. Two Bengali strings therefore both normalised to the empty
string and were scored as nought per cent similar regardless of what they
actually said. Bengali vowel signs are Unicode combining marks and are not
covered by the usual word-character classes, so the bug survived casual
inspection. It was fixed by normalising on Unicode category rather than on a
Latin character class. The general point is larger than the bug: tooling built
on Latin-script assumptions does not fail loudly on a Bengali interface, it
fails silently and reports a plausible number.
\emph{Who is affected:} users and testers of software in Bengali and other
non-Latin scripts. \emph{How the design responds:} script-agnostic
normalisation wherever user-visible text is compared.

\paragraph{Legal considerations.} These are analysed in full in
\Cref{sec:standards}. In summary: dedicated disposable test accounts rather
than production user data, no ingestion of user records into the knowledge
graph, and destructive or communicative actions blocked at the configuration
layer rather than discouraged in prose.

\subsection*{Negative impacts, and what is done about them}

\begin{enumerate}[label=\textbf{N\arabic*.},leftmargin=3.0em]
  \item \textbf{Effect on quality-assurance employment.} A tool whose stated
    purpose is to displace manual effort has a labour effect, and it would be
    dishonest to present it otherwise. The evidence available suggests
    augmentation rather than replacement so far: the closest field deployment in
    the literature reports a saving of 1.5 person-days per regression round on a
    single business line~\cite{auitestagent}, which redirects effort rather than
    removing a role. The design choice that matters here is that the output is a
    natural-language test case a person can read, review and take ownership of,
    not an opaque verdict, so the human role shifts toward judgement rather than
    execution. At larger scale the effect on entry-level testing work is real
    and lies outside what this project can control.
  \item \textbf{Automation bias.} The most likely way this system causes harm is
    not by failing, but by being believed. An organisation that treats an
    automated verdict as authoritative will ship the defects the agent missed.
    \emph{Response:} a reported defect is explicitly scoped as a candidate for
    human confirmation and never as a confirmed fault (\Cref{sec:objectives});
    the generated report states its own degradations rather than hiding them
    (NFR-D2); and the attribution taxonomy makes it structurally impossible for
    the agent's own failure to be filed as an application defect (NFR-D3).
  \item \textbf{Acting on a production system with live users.} The evaluation
    target is in production. An exploratory agent that decides for itself what
    to do is, by construction, capable of doing something irreversible.
    \emph{Response:} the safety boundary is machine-enforced (DR-07) through
    blocked control texts and blocked URL patterns, out-of-scope requirements are
    withheld from the planner's citable list, and accounts used for testing are
    disposable.
  \item \textbf{Prompt injection.} The application's own on-screen text is fed
    to a model that also decides the next action, so content under an attacker's
    control could in principle direct the agent.
    \emph{Response:} every prompt that carries observed application content
    labels it explicitly as untrusted data that must never be read as
    instructions, in line with the OWASP guidance for language-model
    applications~\cite{owaspllm}.
  \item \textbf{False reports wasting developer time.} Each incorrectly reported
    defect costs investigation effort and erodes trust in the tool, which is the
    conflict recorded in \Cref{tab:stakeholders}.
    \emph{Response:} the evidence-integrity requirements above. The honest
    limitation is that the false-positive rate has not been measured, because
    doing so requires the seeded-defect experiment described in
    \Cref{sec:conclusion}.
\end{enumerate}

\section{Environmental Impact and Life Cycle Sustainability}
\label{sec:sustainability}

\subsection*{Development}

The development footprint is five developer workstations, one physical Android
handset, emulator instances on those workstations, and a local graph database.
No hardware was manufactured for the project; the handset is an existing
personal device. The decision that matters most environmentally was taken for
other reasons: model training and fine-tuning were placed out of scope
(\Cref{sec:objectives}), and the project uses hosted inference through an API
instead. Training is the dominant energy cost in most machine-learning work, so
excluding it means this project's footprint is an inference footprint only.

\subsection*{Operation}

The recurring environmental cost is hosted model inference, and it scales with
the number of tokens processed rather than with wall-clock time. Three
properties of the design reduce that volume. All three were adopted to control
money rather than energy, but the mechanism is the same.

\begin{itemize}[leftmargin=1.6em]
  \item Roughly nine tenths of the work in each execution happens on the device
    or in the browser rather than in a model, so most of a campaign consumes no
    inference at all.
  \item The executor and investigator roles run on smaller, cheaper model tiers,
    adopted once the quality cost of the downgrade had been measured.
  \item Context is budgeted rather than accumulated. Findings are atomic and
    deduplicated on write (DR-06), so a discovery made in round one is carried
    into round twenty as a single line with an observation count instead of
    being restated in full each time. Before that change, one test case's
    prose occupied the majority of the next generation prompt.
\end{itemize}

The literature supports token volume as the correct lever. AutoDroid halves its
prompt length by merging functionally equivalent interface nodes and reports
the query cost falling by the same proportion~\cite{autodroid}; AppAgentX
reduces tokens per task from 9.26k to 4.94k by replacing repeated reasoning
with stored shortcuts~\cite{appagentx}. Against these, the per-application
figures for a two-hour exploratory campaign reported by
DroidAgent~\cite{droidagent} indicate the scale of inference that an
unoptimised design consumes.

Device-side energy is smaller but not zero: a handset driven continuously for
the length of a campaign, plus a graph database left running between campaigns.
The honest limitation is that this project measured \emph{monetary} cost, not
energy. Token count is a reasonable proxy, but converting it to energy or
carbon requires the provider's hardware and grid mix, which are not disclosed.
No carbon figure is therefore claimed anywhere in this report.

\subsection*{Data life cycle, which is the real issue}

The knowledge graph accumulates screenshots, accessibility dumps and execution
traces of a production application, gathered through accounts that can sign in
to it. That is sensitive local test evidence with a retention question
attached, not neutral telemetry, and three concrete problems follow.

\begin{enumerate}[label=\textbf{E\arabic*.},leftmargin=3.0em]
  \item \textbf{Decay is not deletion.} A knowledge half-life (default 90 days)
    downweights stale knowledge in retrieval so the agent adapts when an
    application changes. It keeps the graph \emph{relevant}; it does nothing to
    bound how large it grows or how long the underlying evidence persists.
    \emph{Mitigation:} a retention policy that actually deletes, applied to
    screenshots and raw traces separately from the derived findings that give
    them value.
  \item \textbf{Graph snapshots are large, and permanent if mishandled.} A full export of the knowledge graph runs to tens of megabytes and carries the same evidence the live graph does. Version control is the wrong home for such a snapshot, because anything committed to a repository is effectively permanent in its history. \emph{Mitigation:} snapshots are held outside version control, under the same controls as the test credentials and with a stated retention period.
  \item \textbf{Screenshots are the bulk and the most sensitive part.}
    \emph{Mitigation:} store them outside the repository with an explicit
    lifetime, and keep the derived caption and structural signature, which is
    what retrieval actually uses, rather than the image, once that lifetime
    expires.
\end{enumerate}

\subsection*{Retirement}

The system is four processes, one database and no bespoke hardware, so
decommissioning it is deleting a database, revoking API keys and disabling the
test accounts. Nothing becomes electronic waste. The obligation that outlives
the software is the data described above, and it is the part that a retirement
plan has to name explicitly rather than assume will disappear with the code.

\section{Ethics}
\label{sec:ethics}

\subsection{Equity and Inclusivity}
\label{sec:equity}

Three groups could be excluded by a system of this kind, and each was
considered.

\emph{Users of software written in non-Latin scripts.} The normalisation defect
described in \Cref{sec:society} is the concrete instance: a component that
silently scored every pair of Bengali strings as completely dissimilar. It is
representative of a wider pattern in which tooling developed against English
interfaces degrades invisibly elsewhere, and the response was to normalise on
Unicode category so that any script is handled on the same footing.

\emph{Users of assistive technology.} Covered in \Cref{sec:society}: because
the agent is confined to what the accessibility layer exposes, an application
that is opaque to it is generally also opaque to a screen reader, and the
agent's failures make that visible rather than hiding it.

\emph{Organisations without the budget for quality assurance.} The measured
unit cost and the configuration-driven target system are what put this within
reach of a team that cannot fund a dedicated testing function.

The limitation to state honestly is that the agent is only as inclusive as the
interface it is given. An application built with a toolkit that exposes no
stable identifiers and no text through the accessibility tree, which is
precisely constraint C4, is one the agent tests poorly. That is a real
exclusion and it falls on exactly the applications that are already least
accessible.

\subsection{Accountability and Personal Responsibility}
\label{sec:accountability}

Responsibility for a system that issues correctness judgements has to be
located deliberately, because the failure mode is a machine-authored claim that
a human then acts on.

At the system level the position taken throughout this project is that
\textbf{the agent is not accountable for its verdicts; the people operating it
are}. Three mechanisms make that position operable rather than rhetorical. A
reported defect is scoped as a candidate requiring human confirmation
(\Cref{sec:objectives}). The three-way fault attribution taxonomy (NFR-D3) has
exactly one definition in the source tree, so an agent limitation cannot be
reported as an application defect through drift between components. And the
cross-process degradation register (NFR-D2) records every fallback taken, on
the principle that a system that hides its own degradation is worse than one
that fails loudly.

Two decisions during the work illustrate the same principle applied to the
team's own conduct. The team put its own codebase through a structured internal
review against an explicit severity rubric, rather than assuming that code it
had written was code it had checked, and scheduled the resulting work
accordingly. And the evaluation reports autonomy and coverage without
adjustment, including where those numbers are unflattering, because a testing
tool that presents its own results selectively has disqualified itself from the
job.

At the individual level, responsibility follows the subsystem allocation in
\Cref{sec:participation} rather than being held collectively, because a
component owned by everyone is answered for by no one. Each member named in
\Cref{tab:contrib} is the person answerable for the behaviour of the subsystem
in that row: the execution and integration lead for what the device and browser
executors do to an application under test and for the enforcement of the
irreversible-action boundary at the point of execution; the planner and
retrieval owner for what enters a prompt and for the trajectory evaluator's
judgements; the architecture owner for the graph schema, the single-writer
discipline and the fault-attribution definition; the operator-surface owner for
what the dashboard and the exported report assert to a reader who was not
present at the run; and the review and grounding owner for the accuracy of the
literature claims and of the generated narrative in that report.

Collective responsibility is retained for exactly two things, and deliberately:
the decision to report a defect as a candidate rather than a verdict, and the
decision to publish autonomy and coverage unadjusted. Both are positions of the
team as a whole, and neither may be revised by the owner of a subsystem acting
alone.

\subsection{Proper Use of Intellectual Property}
\label{sec:ip}

\paragraph{Prior work and literature.} Every empirical claim taken from the
research literature in \Cref{ch:research} is attributed to the paper it came
from, and the figures were read from those papers directly rather than from
secondary summaries. Where a work was surveyed rather than read in full, this
is stated at the point of use (\Cref{subsec:additional}) instead of being
blurred.

\paragraph{Third-party software.} The system is assembled from open-source
components and hosted services. The principal Python dependencies are FastAPI
and Uvicorn (HTTP services), Pydantic (schemas), the Neo4j driver, LangChain
and LangGraph (agent orchestration), Playwright (browser execution), mobilerun
(Android device execution), fastembed and NumPy (embeddings), markitdown, pypdf
and fpdf2 (document handling), structlog and psutil (operations), and an
OpenRouter client for hosted inference. The dashboard uses React built with
Vite. Neo4j itself is the graph database. Model inference is obtained
commercially through OpenRouter and is governed by that provider's terms of
service.

\Cref{tab:licences} lists the principal third-party components with the licence
under which each is used. Two of them carry copyleft terms and therefore needed
an explicit compatibility judgement rather than an assumption.

\begingroup
\small
\begin{xltabular}{\textwidth}{@{}>{\hsize=0.70\hsize}Y>{\hsize=0.86\hsize}Y>{\hsize=0.44\hsize}Y@{}}
  \caption[Third-party components and their licences]{Principal third-party components and their licences, read from the metadata of the versions actually installed rather than from documentation}
  \label{tab:licences}\\
  \toprule
  Component & Role in the system & Licence \\
  \midrule
  \endfirsthead
  \caption[]{Principal third-party components and their licences \emph{(continued)}}\\
  \toprule
  Component & Role in the system & Licence \\
  \midrule
  \endhead
  \midrule
  \multicolumn{3}{r@{}}{\footnotesize\itshape continued on the next page}\\
  \endfoot
  \bottomrule
  \endlastfoot

  FastAPI, Uvicorn, Pydantic & Service framework, ASGI server, schema validation & MIT, BSD-3, MIT \\
  \addlinespace
  Neo4j Community Edition & Knowledge-graph database (self-hosted server) & GPL-3.0 \\
  Neo4j Python driver & Client access over the Bolt protocol & Apache-2.0 \\
  \addlinespace
  LangChain, LangGraph & Agent orchestration & MIT \\
  Playwright (Python) & Browser execution & Apache-2.0 \\
  mobilerun 0.6.11 & Android device execution through the accessibility service & MIT \\
  \addlinespace
  fastembed, onnxruntime, NumPy & Local embeddings, their inference runtime, and numerical support & Apache-2.0, MIT, BSD-3 \\
  markitdown, pypdf & Format-agnostic document ingestion & MIT, BSD-3 \\
  fpdf2 & User-guide PDF generation (\modpath{scripts/make\_user\_guide.py}) & LGPL-3.0 \\
  \addlinespace
  structlog, psutil, requests, python-dotenv & Logging, process inspection, HTTP, configuration & Apache-2.0/MIT, BSD-3, Apache-2.0, BSD-3 \\
  \addlinespace
  React, React DOM, Vite, \modpath{@vitejs/plugin-react}, \modpath{vite-plugin-singlefile} & Operator dashboard, built to a single self-contained file & MIT \\
  \addlinespace
  Hosted model inference (OpenRouter; Gemini) & Commercial services, not distributed with the system & Provider terms of service \\
  \bottomrule
\end{xltabular}
\endgroup

\paragraph{Compatibility.} The database server is distributed under GPL-3.0,
but the system does not link it: it communicates with a separate server process
over the Bolt network protocol using an Apache-2.0 client driver, which is the
arrangement under which the server's copyleft does not extend to this
project's own source. The PDF generator is LGPL-3.0 and is imported as a
library, which the licence permits without relicensing the importing work,
subject to the recipient being able to substitute a modified version of the
library. Everything else is permissive. The practical consequence is a
constraint on \emph{redistribution} rather than on use: a hosted or internal
deployment as described in \Cref{sec:commercialisation} is unaffected, whereas
shipping a single bundled artefact containing the database server would have to
honour GPL-3.0.

\paragraph{The team's own code.} The project is released under the Apache
License 2.0, and the licence text is in the repository root. The choice was
made on two grounds. It is compatible with every component in
\Cref{tab:licences}, including the GPL-3.0 database server the system speaks to
over the network and the LGPL-3.0 library it imports, so the release creates no
obligation the project cannot meet. And unlike the shorter permissive licences
it carries an explicit patent grant, which matters for work carried out with an
industrial partner: it makes the terms on which the partner and any later
adopter may use the result explicit rather than implied. This is the licence
for the code the team wrote. It does not extend to the partner's requirements
document, which is not the team's to license, and it does not alter any term of
the engagement under which the work was commissioned.

\paragraph{The partner's material.} The technical requirements document
ETA-TR-2026-001 is not public. It is cited for traceability, and requirement
identifiers and titles are reproduced where they are needed to show what was
delivered, but no confidential content beyond that appears in this report.

\paragraph{Generative AI.} Generative AI tools were used during this project,
and the Candidates' Declaration requires that use to be disclosed. The team
used an agentic coding assistant, under a repository-level instruction file the
team wrote and maintained, in four areas: drafting and refactoring
implementation code within a design the team had already fixed; generating and
maintaining internal documentation; assisting with the prose of this report;
and producing the connecting text in the system's own exported run report,
where every number is computed from the knowledge graph.

Judgement was not delegated. The four-process architecture, the decision that
the components communicate only through a persistent graph, the requirements
analysis and its traceability, the fault-attribution taxonomy, the choice of
what to measure and the interpretation of the results are the team's own.
Generated code was read, executed and tested before it was kept, and factual
statements in this report were checked against a primary source rather than
accepted from a model: configuration against the source files, chronology and
contribution against the version-control record, costs against the provider's
billing record, licences against installed package metadata, and the figures
attributed to other systems in \Cref{ch:research} against the papers
themselves.

\subsection{Professionalism and Ethical Codes}
\label{sec:codes}

The ACM Code of Ethics and the IEEE Code of Ethics both place an obligation on
the engineer to avoid harm, to be honest about the limitations of their work,
and to accept responsibility for professional judgements. The IEB code makes
the same demand of the engineer's duty to the public. Three decisions in this
project were settled by those obligations rather than by convenience, and each
is documented at the point where it was taken.

\emph{Honesty about what the work establishes.} The review in
\Cref{subsec:measurement} concludes that this project cannot report precision or
recall for defect detection without seeding known faults, and says so, rather
than substituting a coverage number that would look like a result. The
corresponding ACM obligation is to be honest and trustworthy, including about
the limits of one's own competence and evidence.

\emph{Avoiding harm to a third party's users.} The evaluation target is a
production system with real users. Registration, one-time-password
verification, payment and account deletion were placed out of scope and blocked
in configuration rather than merely discouraged (DR-07), on the straightforward
reading that an engineer may not risk another party's users to make their own
evaluation more complete.

\emph{Applying to our own system the standard we ask of others.} The test account's secret is passed only to the component that has to type it, and never to the planner, which receives the account role alone (NFR-D6). The planner would produce marginally better-grounded tests with the full credential in its context, so the decision cost a little capability to uphold a principle that would have been invisible had it gone the other way. An engineer's obligations do not lapse because the system under discussion is their own.
